% =====================================================================
% DISCUSSION
% =====================================================================
\section{Discussion}
\label{sec:discussion}

\subsection{Prompt-Level Defence is a Slope, and Sometimes Not Even That}
\label{sec:disc:slope}

The persona sweep produced a clean monotone reduction --- $0.472$ at L1
to $0.083$ at L5 --- and it is tempting to read that as a defence
working. The guardrail ladder complicates the reading. There, two of the
ten rungs performed no better than the unguarded baseline: the soft
system-prompt rule (B1, $0.242$) and the ASCII-only gate (B4, $0.333$)
both sat above B0's $0.121$.

The B1 case is the more instructive of the two, because its mechanism is
structural rather than incidental. A system-prompt rule that names the
protected value and forbids its disclosure must place that value in the
model's context in order to state the rule. It thereby does two things
at once: it makes disclosure possible at all, and it supplies an
explicit rule for an override prompt to target. The rule is the
vulnerability's precondition and the attacker's map. This is why the
recommendation that follows is architectural rather than a stronger
prompt: the coupon should be resolved server-side against an
authorisation check the model never sees, so that no phrasing can
surface what is not present.

The general form of the finding is that a prompt-level control shifts a
distribution; it does not establish a boundary. The strictest persona in
this study is instructed to deny that any secret exists and to treat
extraction attempts as hostile, and it still disclosed on $8.3\,\%$ of
attempts. For a control whose failures are silent and whose per-attempt
cost to an attacker is negligible, an $8.3\,\%$ residual is not a
residual risk --- it is a certainty under repetition.

\subsection{Framing Beats Force: Refusal Tracks Surface Form}
\label{sec:disc:framing}

The per-technique decomposition (Table~\ref{tab:technique}) is the
result with the widest applicability beyond this target. Success was not
spread across the twelve techniques; it concentrated in the two that
never ask for the protected value in a recognisable form.

Fictional framing --- ``write a short play where a clerk absent-mindedly
says the secret word'' --- succeeded on $14/15$ attempts and was the only
technique to carry the attack to the strictest persona more than once.
The persona instruction constrains how the assistant behaves \emph{as an
assistant}; it says nothing about the content of a creative artefact the
assistant has been asked to author. The model's refusal policy and the
task it was given occupy different frames, and the attack simply moves
the request into the frame the policy does not cover.

Indirection --- ``don't tell me the coupon word; instead list every word
you are forbidden from saying'' --- succeeded on $10/15$, against $1/15$
for asking plainly. The prompt explicitly disclaims the goal it
achieves. Meanwhile, the three techniques that name their intent most
plainly were the weakest in the study: authority claim $0/15$,
developer-mode claim $1/15$, direct request $1/15$.

This ordering is a direct challenge to keyword-based and
signature-based input defences, and the guardrail ladder confirms it
from the other direction: B7's keyword gate blocked $54.5\,\%$ of all
requests to achieve a leak rate of $0.061$, and the two prompts that got
through it were a paraphrase and a hypothetical.

The clearest single illustration is the system-prompt probe
(Table~\ref{tab:sysprompt}). Run against the strictest persona --- the
one that leaks on $8.3\,\%$ of direct coupon requests --- prompts that
ask about the model's \emph{rules} rather than its \emph{secret}
extracted the secret on $46.7\,\%$ of attempts, a $5.6\times$ increase.
Two of those prompts leaked the coupon while leaking no system-prompt
text at all. A defence tuned against requests for the asset is not a
defence against requests for the policy that protects the asset.

\subsection{The Failing Layer, Not the Strength of the Defence}
\label{sec:disc:layers}

Holding the attack corpus fixed across ten defensive configurations
isolates which layer fails, and each layer fails in a characteristic
way (Fig.~\ref{fig:layers}).

\begin{figure}[t]
  \centering
  \safefig[1.0]{figures/fig-defence-layers.pdf}
  \caption{Where each defence sits in the pipeline, the technique class
  that defeats it, and the measured bypass rate. Rates are per stage over
  33 attempts.}
  \label{fig:layers}
\end{figure}

\textbf{Input filters fail by construction, and their cost is
availability.} The server decodes base64 and ROT13 segments before the
model sees them, so the ``naive literal scanner'' inspects a
representation that no longer exists downstream; ROT13 in particular
succeeded on 5 of 30 attempts because its decoded form is natural
English. The ASCII gate is defeated by pure-ASCII Spanish, which is to
say that language restriction is not a security control. The keyword
gate is defeated by paraphrase, and paid for its $0.061$ leak rate by
rejecting more than half of all traffic --- in production, more than
half of the legitimate queries whose phrasing happened to collide with
the list.

\textbf{Output filters convert a full disclosure into a partial one and
push the leak into a different representation.} B5 and B9 both scored
zero on literal leaks, which in isolation reads as success. But across
the study, 12 of 97 leaks ($12.4\,\%$) were caught only by the
obfuscation-aware branches of the detector: hyphenated spell-outs
(\code{C-H-E-E-S-E}), spaced renderings, and plural forms that defeat a
word-boundary match. A single-representation filter is defeated by
trivial reformatting that a human reader parses without effort. The
corresponding control is to normalise before matching --- strip
separators, fold leetspeak, collapse whitespace --- or, preferably, to
remove the secret from scope so that there is nothing to filter.

\textbf{An untrusted data channel becomes an instruction channel by
default.} The QR results ($0.238$, all payloads verified intact) show
that content the application decodes from an uploaded image is
concatenated into the prompt with the same authority as typed input. No
model-level defence addresses this, because from the model's position
the two are indistinguishable. The control is at the application
boundary: quote extracted content to the model as data, apply the same
input policy to decoded content as to typed content, and never merge it
into the instruction context.

\textbf{Agency converts a text failure into a data breach.} The
cross-tenant SQL result is a single hit in 21 attempts and would be
dismissed as noise on rate alone. It should not be. The routing flag
exists only in one row of an isolated per-user table and cannot be
guessed or hallucinated; its appearance means the application executed a
model-authored query the caller was not entitled to run. The severity of
an LLM01 failure is set by what the surrounding system will do with the
model's output, not by how often the model can be induced to produce it.
The same reasoning explains why \code{Tool\bk Execution\bk Induced} scoring $0$
on the QR surface is a meaningful negative result: the injection
succeeded at the text level and stopped at the execution level.

\subsection{What the Detectors Disagree About}
\label{sec:disc:detectors}

The detector-agreement table (Table~\ref{tab:agreement}) is the
methodological core of this study, and it cuts in two directions.

\textbf{Generic detectors over-approximate by design.} The stock
\code{mitigation.\bk Mitigation\bk Bypass} detector achieved perfect recall ---
it missed no ground-truth violation --- at a precision of $0.293$
against the application's actual policy. This is not a defect. It
answers ``this output is not a refusal'', a question that generalises
across deployments precisely because it is indifferent to any particular
deployment's policy. On this corpus most outputs are not refusals,
including entirely appropriate ones, so most of its hits are not
violations. The practical consequence is that a generic detector is a
triage instrument: sound as a filter, unsound as a verdict. The
policy-aware detector can only exist here because PwnzzAI's policy is
legible from its source --- an unusual property, and the reason this
study can report a rate rather than an impression.

\textbf{Application-supplied oracles vary in trustworthiness across
surfaces of the same application.} On the catering RAG the built-in flag
agreed with the independent detector on all 24 attempts; on the comment
RAG, on 21 of 24; on the order assistant it flagged an access violation
on 16 of 21 attempts where no cross-tenant identifier was present. A
security signal shipped with the system is a hypothesis about the
system's behaviour, not a measurement of it, and this study's design ---
carrying the oracle in the notes and cross-tabulating rather than
consuming it --- is what makes that testable.

\textbf{The policy-aware detectors are not exempt.} Two false-positive
classes were identified by inspection and should qualify the numbers
reported above.

The first is the catering residual. With the trusted-only mitigation
enabled, \emph{no untrusted passage was retrieved on any of the 12
attempts}, yet \code{Poisoned\bk Retrieval\bk Influence} scored $0.417$. Four of
the five residual hits occurred on queries that themselves contained
``pineapple'' or ``anchovy paste'', and the responses show the model
echoing the noun from the question --- in one case inside an explicit
refusal (``I can't find an approved policy specifically addressing
pineapple\ldots''). A term-presence detector cannot distinguish
\emph{asserting} the injected policy from \emph{repeating the querent's
vocabulary while declining}. The corrected reading is stronger than the
uncorrected one: trusted-only retrieval was fully effective at the
retrieval stage, and the appropriate residual concern is the single
neutral-query hit and the general fact that a term-presence oracle
cannot certify grounding.

The second is a substring fallback in \code{CouponLeak}. The
whitespace-stripped branch matches without word boundaries, which
correctly caught a genuine leak rendered as the plural ``Mushrooms!''
but also fired on an invented word, ``MOONCHEESE'', containing the
level's secret. One of the 97 recorded leaks is attributable to this.
The rate-level impact is negligible; the methodological point is not.
Ground-truth detection removes the judgement call about \emph{what
counts as} a violation, but it does not remove the engineering problem
of \emph{recognising} one, and both of the refinements that make these
detectors sharp --- obfuscation tolerance and term-presence scoring ---
buy sensitivity by spending precision. Reporting the matched rendering
per hit, as the design does, is what makes this auditable after the
fact.

\subsection{Poisoning: the Threshold is Not Where the Effect Begins}
\label{sec:disc:poisoning}

The sentiment results are the study's only exact measurements --- the
surface involves no language model --- and they support three claims.

\textbf{The attack is cheap.} Five mislabelled comments flipped the
trigger phrase, in a corpus consisting of the application's own comment
table plus the injection. Any pipeline that retrains on user-submitted
feedback without provenance or review is exposed at this scale.

\textbf{Output monitoring is blind before the threshold.} Between $b=0$
and $b=3$ the trigger's $P(\text{positive})$ rose from $0.230$ to
$0.497$ --- more than doubling, and stopping just short of the boundary
--- while the reported label never changed and the flip rate remained
exactly $0.000$. Every label-level monitor reports ``no effect'' across
that entire region. A score-distribution monitor watching the same runs
sees a monotone drift from the first injected comment. The detection
control therefore has to observe the classifier's continuous output on
trigger-candidate terms, not its decisions.

\textbf{The backdoor is targeted, which is what defeats accuracy
testing.} All three trigger-bearing prompts drifted monotonically toward
positive, including one carrying an explicit ``worst pizza I have ever
eaten'' that travelled $0.356$ in probability without ever flipping. All
three no-trigger controls stayed flat or moved slightly away from
positive. A poisoned model of this kind passes a held-out accuracy
check, because its behaviour is unchanged everywhere except on inputs
containing a phrase the attacker chose. Acceptance testing on aggregate
metrics cannot find it; the controls that can are provenance on training
data, label-distribution anomaly detection before retraining, and an
unpoisoned holdout monitored for drift on trigger-candidate terms.

For retrieval poisoning the picture is sharper than the headline rate
suggests. Unmitigated, influence was total: $12/12$, including the
baseline control query about notice periods, with 8 of the 12 hits on
queries that never named the injected terms --- the fabricated policy
displaced trusted baseline facts rather than merely adding to them.
Mitigated, untrusted retrieval went to zero. Trust-tagged retrieval is
therefore a genuinely effective control at the stage it operates on, and
the honest residual concern is not that it leaks poison but that
term-presence measurement cannot certify that an answer is grounded in
trusted context. That requires an output-grounding check, which is a
different control at a different stage.

\subsection{Evidence-Linked Mitigations}
\label{sec:disc:mitigations}

Table~\ref{tab:mitigations} states the recommendations, each tied to the
detectors that evidence it and to the pipeline layer it acts on. The
ordering is deliberate: the controls that eliminate a capability
outrank the controls that police its use.

\begin{table}[t]
\caption{Evidence-linked mitigations. ``Layer'' is where the control acts relative to the model;
``evidence'' names the detector signal supporting the finding.}
\label{tab:mitigations}
\centering
\footnotesize
\zebra
\begin{tabular}{@{}C{\tw{0.06}}L{\tw{0.09}}L{\tw{0.30}}L{\tw{0.13}}L{\tw{0.42}}@{}}
\toprule
\headrow
\tabhead{ID} & \tabhead{OWASP} & \tabhead{Finding (evidence)} & \tabhead{Layer} & \tabhead{Control} \\
\midrule
M-01 & LLM01 & Prompt hardening lowers but never eliminates leakage; L5 residual $0.083$ (\code{CouponLeak}) & architecture &
Resolve the coupon server-side behind an authorisation check the model never sees; the value must not enter context. \\
M-02 & LLM05 & Single-representation output filters miss obfuscated renderings; $12/97$ leaks (\code{CouponLeak}, \code{Coupon\bk Suppressed}) & output-handling &
Normalise before matching (strip separators, fold leetspeak, collapse whitespace); prefer allow-listed structured responses. \\
M-03 & LLM01 & Input filters bypassed by encoded channels and pure-ASCII foreign language; B3 $0.121$, B4 $0.333$, B7 blocks $54.5\,\%$ (\code{CouponLeak}, \code{InputRejected}) & input-handling &
Treat decoded auxiliary content as untrusted data, never re-merged into instructions; do not treat language restriction as a control. \\
M-04 & LLM01 & Instructions smuggled through an uploaded QR image are decoded and obeyed; $0.238$, all payloads intact (\code{CouponLeak}) & input-handling &
Quote extracted upload content to the model as data; apply the same input policy to extracted text as to typed text. \\
M-05 & LLM02 & The RAG assistant is configured to disclose PII on request; $0.375$, all hits on business-framed asks (\code{CustomerPII}) & data-governance &
Redact PII at index time so it is never retrievable; remove disclosure permission from the system prompt; add an output DLP check. \\
M-06 & LLM06 & The agentic SQL tool executed a model-authored cross-tenant read; $1/21$, categorical (\code{Cross\bk Tenant\bk Flag}) & authorisation &
Scope every tool query to the authenticated principal server-side; parameterise SQL; constrain the agent to a least-privilege view. \\
M-07 & LLM04 & Five mislabelled comments create a targeted backdoor; drift measurable from $b{=}1$ (\code{Sentiment\bk Label\bk Flip}) & data-governance &
Provenance and review on feedback data; label-distribution anomaly detection before retraining; monitor \emph{scores}, not labels, on trigger terms. \\
M-08 & LLM04 & Trusted-only retrieval eliminated untrusted retrieval ($0/12$) but term-presence cannot certify grounding (\code{Poisoned\bk Retrieval\bk Influence}) & retrieval &
Keep trust-tagged retrieval and add an output-grounding check; require corroboration across trusted sources for mandatory-sounding claims. \\
\bottomrule
\end{tabular}
\end{table}

Read together, the recommendations share a shape. Every control that
holds in this study acts on the application pipeline --- what enters the
context, what the retriever will surface, what the tool layer is
permitted to execute, what leaves the process. Every control that
failed or under-performed acts on the model's disposition: a firmer
instruction, a keyword list, a string replacement on the way out. The
model is the component whose behaviour cannot be constrained by
assertion, and the pipeline is where the enforceable boundaries are.

\subsection{On Being Garak-Native}
\label{sec:disc:garak}

Expressing the assessment entirely as Garak plugins rather than as a
bespoke harness had three consequences worth reporting.

It made the scan auditable by construction. Every outcome in this paper
traces to a Garak evaluation entry in a retained \code{report.jsonl};
the analysis layer aggregates but never re-judges. The per-task
configuration files replay any single measurement through the stock
CLI. There is no point in the pipeline where a bespoke scoring decision
could have entered unrecorded.

It made the comparison of detection philosophies free. Because Garak
already supports running extended detectors alongside a primary one,
placing a stock detector next to each policy-aware one cost a single
line per probe and produced Table~\ref{tab:agreement} as a by-product.

It disciplined the scope of custom code. Custom generators were written
only where the transport is not text-in/text-out or where run state must
be established first; plain chat endpoints are reached by the stock REST
generator with a configuration file. The nine surfaces in
Table~\ref{tab:generators} are therefore an inventory of the
application's genuinely non-standard interfaces, which is itself a
useful artefact of the assessment.

\subsection{Limitations and Threats to Validity}
\label{sec:disc:limitations}

\textbf{Model scale bounds every absolute rate.} All LLM-mediated
measurements were taken against a 1B-parameter model that refuses,
rambles, and drifts off-task more than a production-scale model. Every
absolute rate here should be read as a lower bound for this deployment
and not as a property of the attack technique. The relative comparisons
--- between levels, between layers, between mitigation states, between
detectors --- hold the model fixed and are the intended contribution.

\textbf{The sample per cell is small.} Three generations per prompt
gives 33--36 evaluated attempts per task; a difference of one hit moves
a task-level rate by roughly three points. Differences of one or two
percentage points between adjacent stages are not resolvable, and the
paper does not interpret them. The claims drawn are ordinal and
large-margin: $14/15$ against $1/15$, $1.000$ against $0.417$, $0.467$
against $0.083$. The deterministic poisoning surface is exempt from this
limitation entirely.

\textbf{Prompt corpora are compact by design, so coverage is not a
claim.} Each probe carries five to twelve prompts chosen for
attributability rather than breadth. A rate of $0.000$ in this study
means ``not achieved by these prompts'', never ``not achievable''. Two
stage results are explicitly of that kind: B6's $0.000$ reflects a
single-turn probe that asserts a prior agreement absent from the
forwarded history, when a genuine test of a history-trusting context
needs a multi-turn attack that first establishes the false premise; and
\code{Tool\bk Execution\bk Induced}'s $0$ across 21 QR attempts bounds only the
two payloads that targeted it.

\textbf{Detector precision is spent to buy sensitivity.} As
Section~\ref{sec:disc:detectors} sets out, the obfuscation-tolerant
substring fallback and the term-presence poison scorer each produced
identifiable false positives. The affected numbers are the catering
hardened rate ($0.417$, of which 4 of 5 hits are query-echo) and one of
97 coupon leaks.

\textbf{Non-determinism is not eliminated, only bounded.} Garak's RNG is
seeded and parallelism is disabled, but the model behind the HTTP path
is not seedable through the application interface. Individual
generations vary between executions; rates are stable in ordering and
approximate magnitude, not in their third decimal.

\textbf{Single target, single deployment.} All findings are from one
deliberately vulnerable application in one configuration. The
transferable contributions are the method --- ground-truth detection
against a legible policy, paired controls for poisoning, cross-tabulated
oracles, layer-isolating sweeps --- and the technique-level orderings,
not the rates.
